\section{Reliability, Long Horizons and Cost}\label{sec:reliability}

Before verification comes a prior question for any unattended run: how often, how consistently and at what price the agent finishes. The evidence gives four answers. Success falls steeply with task length; run-to-run variation is large enough to swamp most reported differences; humans still steer most production work; and cost per task is large, variable and loosely tied to accuracy.

\subsection{Error Compounding and Horizon Length}

The headline horizon overstates the reliable one. METR estimates that o3 completes tasks of about 110 human-minutes half of the time, a horizon that has doubled roughly every seven months since 2019, but finds that ``80\% time horizons are 4-6x shorter'' \citep{F1}. A second METR benchmark finds agents succeeding 70--80\% of the time on tasks that take humans under an hour and under 20\% on tasks over four hours \citep{F38}. The mechanism is compounding. When the plan and the knowledge are supplied, the best model in one study falls below 50\% accuracy within 15 turns; injected errors degrade later turns further, and larger models do not escape this ``self-conditioning'', although thinking models do \citep{F3}. Early mistakes cascade \citep{F5}: across 2,518 annotated trajectories, agents recover after a first mistake in only 30.5\% of runs and never detect it in 38.5\%, and some runs scored as solved ``fabricate success'' \citep{F25}. AstaBench states the arithmetic: with about ten steps at 70\% each, completing all of them has probability $0.7^{10} \approx 3\%$, and the observed maximum is 5\% even though step scores reach about 70\% \citep[p.~10]{S19}. The same arithmetic cuts the other way: small gains in per-step accuracy compound into exponential gains in the length of tasks completed \citep{F3}. Both readings imply that a horizon claim should state the success rate it assumes.

\subsection{Variance: Why One Run Is Not a Result}

A single run per task cannot rank agents. Across 60,000 trajectories on SWE-bench Verified, a single run's pass@1 moved by 2.2 to 6.0 percentage points depending on which run was taken, with standard deviations above 1.5 points even at temperature 0, so ``reported improvements of 2--3 percentage points may reflect evaluation noise'' \citep{F19}. The same study finds gaps of up to 24.9 points between pass@k, the chance that at least one of $k$ attempts succeeds, and pass\textasciicircum k, the chance that all $k$ do. On $\tau$-bench, gpt-4o falls from about 61\% at pass\textasciicircum 1 to about 25\% at pass\textasciicircum 8 on the retail domain \citep{A9}. In 584 runs on one machine-learning task, identical runs of one agent--model pairing ``varied more than the pairings differed'', and separating the agents would take ``tens to more than a hundred runs of each'' \citep{F24}. Across 15 models, capability gains have yielded ``only small improvements in reliability'', and smaller models are often as consistent as larger ones or more so \citep{F2}. Infrastructure adds noise: in a post by Anthropic about its own evaluation setup, container resource limits alone moved Terminal-Bench 2.0 scores by 6 points ($p < 0.01$) \citep{F27}.

\subsection{Humans Still in the Loop}

The available data on production autonomy come mostly from labs measuring their own systems, and they agree that humans still steer most work. Google DeepMind reports that fully delegated pull requests rose from 7.5\% to 11.9\% between April and September 2026, while ``about 88\% of pull requests still require iterative multi-turn human steering or human code authorship'' \citep[sec.~3.4.2]{F30}. OpenAI estimates that autonomous success rose from 23\% to 46\% on two-to-four-hour tasks and from 13\% to 35\% on four-to-eight-hour tasks, yet over half of its successful four-to-eight-hour tasks involved at least one intervention, and it calls its measurement ``preliminary'' \citep{F32}. Anthropic reports interventions per Claude Code session falling from 5.4 to 3.3, from one provider's traffic, as classified by Claude \citep{F14}.

\subsection{Cost}

Cost per task is large, highly variable, and a weak predictor of accuracy. Agentic coding consumes about 1,000 times the tokens of code chat, runs on one task differ by up to 30 times, accuracy ``often peaks at intermediate cost'', and models underestimate their own token use \citep{F6}. In 21,730 rollouts costing about \$40,000, the costliest model run lay on the cost--accuracy frontier in only 1 of 9 benchmarks, and higher reasoning effort gave equal or lower accuracy in 21 of 36 comparisons \citep{F7}. Harness choice alone changed cost by up to 2.6 times on the same model and task \citep{F23},. Since agents of similar accuracy can differ in cost by almost two orders of magnitude \citep{A10}, cost-of-pass, the expected dollars per correct solution \citep{F8}, is a more honest unit than accuracy alone.

\subsection{Practitioner Bounds}

Practitioner tools use several independent caps. The minimal agent used as a SWE-bench baseline stops on separate step, dollar (default \$3), wall-clock and repeated-format-error limits, and records the stop as a named exit status (\repofile{SWE-agent/mini-swe-agent}{src/minisweagent/agents/default.py}); it saves the trajectory after every turn because agents killed at a timeout left none (\repo{SWE-agent/mini-swe-agent}{585}). The UK AI Security Institute's evaluation framework adds time, working-time, message, token, turn and cost limits (\href{https://github.com/UKGovernmentBEIS/inspect_ai/blob/HEAD/docs/setting-limits.qmd}{\texttt{\small UKGovernmentBEIS/inspect\_ai}: \texttt{\small docs/setting-limits.qmd}}). Runs do hang: one user reports that about 20\% of Claude Code review runs never respond, so the GitHub Action times out after six hours (\repo{anthropics/claude-code-action}{674}). But the bounds are instruments that can fail too: a watchdog in one multi-agent orchestrator misidentified healthy sessions and ``killed 22 sessions across all 9 active rigs'' (\repo{gastownhall/gastown}{2707}). For builders, the evidence implies that a long run needs several caps, repeated trials and a supervisor that is itself checked.
